\chapter{Penjumlahan dan Pengurangan}\label{ch:g3:addsub}

Menjumlahkan berarti menggabungkan; mengurangkan berarti mengambil,
atau mencari bagian yang hilang. Bab ini melatih dua cara bersusun
--- dengan simpanan dan pinjamannya yang terkenal --- serta trik
menghitung di kepala yang sering membuat cara bersusun tidak perlu.

\section{Penjumlahan bersusun}

\begin{method}[Menjumlahkan secara bersusun]\label{met:g3:addsub:add}
\begin{enumerate}
  \item Tulislah bilangannya satu di bawah yang lain, satuan di bawah
        satuan, puluhan di bawah puluhan (rata di \emph{kanan});
  \item jumlahkan satuannya; jika \omterm{def:g1:addition:def}{jumlahnya} melewati $9$, tulislah
        angka satuannya lalu \emph{simpanlah} puluhannya ke kolom
        berikutnya (sebagai angka $1$ kecil);
  \item lanjutkan dengan puluhan, ratusan, \dots, sambil menambahkan
        simpanan setiap kali.
\end{enumerate}
\end{method}

\begin{example}\label{ex:g3:addsub:add}
Hitunglah $457 + 368$:
\[
\begin{array}{r}
{\scriptstyle 1}\ \ {\scriptstyle 1}\ \ \phantom{0} \\[-3pt]
4\,5\,7 \\
+\ 3\,6\,8 \\
\hline
8\,2\,5
\end{array}
\]
Satuan: $7 + 8 = 15$: tulis $5$, simpan $1$. Puluhan:
$5 + 6 + 1 = 12$: tulis $2$, simpan $1$. Ratusan: $4 + 3 + 1 = 8$.
Hasilnya: $825$.
\end{example}

\begin{example}[Menjumlahkan di kepala]\label{ex:g3:addsub:mental}
Carilah bilangan yang ramah sebelum memakai cara bersusun:
\begin{itemize}
  \item \emph{buatlah puluhan}: $47 + 8 = 47 + 3 + 5 = 50 + 5 = 55$;
  \item \emph{lompatlah per bagian}: $256 + 130 = 256 + 100 + 30 = 386$;
  \item \emph{tukarlah urutannya}: $17 + 59 + 3 = (17 + 3) + 59 = 79$.
\end{itemize}
\end{example}

\begin{omfigure}
\begin{tikzpicture}[scale=1.05]
  \draw[-Stealth] (-0.2,0) -- (10.4,0);
  \foreach \x/\l in {0.56/{256}, 5.56/{356}, 7.06/{386}}
    { \draw (\x,-0.08) -- (\x,0.08);
      \node[below, font=\small] at (\x,-0.1) {$\l$}; }
  \draw[-Stealth, omDef, thick] (0.56,0.25) .. controls (2.8,1.15) ..
    (5.56,0.25);
  \node[omDef, font=\small] at (3.05,1.05) {$+100$};
  \draw[-Stealth, omThm, thick] (5.56,0.25) .. controls (6.3,0.8) ..
    (7.06,0.25);
  \node[omThm, font=\small] at (6.3,0.85) {$+30$};
\end{tikzpicture}

{\small Menjumlahkan di kepala sebagai lompatan di garis bilangan:
$256 + 130$ dalam dua lompatan, ratusannya dulu, lalu puluhannya.}
\end{omfigure}

\section{Pengurangan bersusun}

\begin{method}[Mengurangkan secara bersusun]\label{met:g3:addsub:sub}
\begin{enumerate}
  \item Tulislah bilangan yang lebih besar di atas, rata di kanan;
  \item kurangkan kolom demi kolom, mulai dari satuan;
  \item jika angka di atas terlalu kecil, \emph{pinjamlah}: ambil
        satu dari tempat berikutnya pada baris atas (nilainya
        sepuluh di kolom yang sedang dikerjakan);
  \item \emph{periksalah}: hasilnya ditambah bilangan pengurang
        harus mengembalikan bilangan di baris atas.
\end{enumerate}
\end{method}

\begin{example}\label{ex:g3:addsub:sub}
Hitunglah $603 - 147$:
\[
\begin{array}{r}
6\,0\,3 \\
-\ 1\,4\,7 \\
\hline
4\,5\,6
\end{array}
\]
Satuan: $3 - 7$ tidak mungkin; pinjamlah lewat tempat berikutnya:
$603$ menjadi $5$ ratusan, $9$ puluhan, $13$ satuan. Lalu
$13 - 7 = 6$; puluhan $9 - 4 = 5$; ratusan $5 - 1 = 4$. Hasilnya:
$456$. Periksa: $456 + 147 = 603$.
\end{example}

\begin{omfigure}
\begin{tikzpicture}[scale=0.82]
  % sebelum: 6 ratusan, 0 puluhan, 3 satuan
  \node[font=\small] at (2.5,2.25) {$603$ sebelum meminjam};
  \foreach \i in {0,...,5}
    \draw[omProp, fill=omProp!25] (0.62*\i,0.7) rectangle
      (0.62*\i+0.44,1.7);
  \node[below, font=\footnotesize] at (1.7,0.55) {$6$ ratusan};
  \node[font=\footnotesize] at (4.35,1.2) {$0$ puluhan};
  \foreach \i in {0,1,2} \fill[omDef] (5.4+0.4*\i,1.2) circle (0.13);
  \node[below, font=\footnotesize] at (5.8,0.55) {$3$ satuan};
  % sesudah: 5 ratusan, 9 puluhan, 13 satuan
  \begin{scope}[yshift=-3.3cm]
  \node[font=\small] at (2.5,2.25) {$603$ sesudah meminjam};
  \foreach \i in {0,...,4}
    \draw[omProp, fill=omProp!25] (0.62*\i,0.7) rectangle
      (0.62*\i+0.44,1.7);
  \node[below, font=\footnotesize] at (1.4,0.55) {$5$ ratusan};
  \foreach \i in {0,...,8}
    \draw[omThm, very thick] (3.5+0.22*\i,0.7) -- (3.5+0.22*\i,1.7);
  \node[below, font=\footnotesize] at (4.4,0.55) {$9$ puluhan};
  \foreach \i in {0,...,6} \fill[omDef] (6.1+0.4*\i,1.45) circle (0.13);
  \foreach \i in {0,...,5} \fill[omDef] (6.1+0.4*\i,0.95) circle (0.13);
  \node[below, font=\footnotesize] at (7.3,0.55) {$13$ satuan};
  \end{scope}
\end{tikzpicture}

{\small Pinjaman pada $603 - 147$ dalam gambar: satu ratusan dibuka
menjadi sepuluh puluhan, dan satu dari puluhan itu menjadi sepuluh
satuan. Banyaknya tetap sama, hanya bajunya yang baru ---
$6\,\vert\,0\,\vert\,3$ menjadi $5\,\vert\,9\,\vert\,13$, dan
sekarang setiap kolom dapat dikurangkan.}
\end{omfigure}

\begin{example}[Pengurangan sebagai bagian yang hilang]\label{ex:g3:addsub:missing}
``Lila punya $35$ kelereng, Tono punya $52$: berapa lebihnya Tono?''
adalah pengurangan $52 - 35$. Di kepala, bilanglah \emph{maju} dari
$35$: $35 \to 40$ adalah $5$, lalu $40 \to 52$ adalah $12$;
seluruhnya $5 + 12 = 17$. Membilang maju sering lebih mudah daripada
meminjam.
\end{example}

\begin{omfigure}
\begin{tikzpicture}[scale=1.05]
  \draw[-Stealth] (-0.2,0) -- (10.4,0);
  \foreach \x/\l in {1.0/{35}, 3.5/{40}, 9.5/{52}}
    { \draw (\x,-0.08) -- (\x,0.08);
      \node[below, font=\small] at (\x,-0.1) {$\l$}; }
  \draw[-Stealth, omDef, thick] (1.0,0.25) .. controls (2.2,0.9) ..
    (3.5,0.25);
  \node[omDef, font=\small] at (2.25,0.92) {$+5$};
  \draw[-Stealth, omThm, thick] (3.5,0.25) .. controls (6.5,1.2) ..
    (9.5,0.25);
  \node[omThm, font=\small] at (6.5,1.1) {$+12$};
\end{tikzpicture}

{\small Menghitung $52 - 35$ dengan membilang maju: mula-mula ke
$40$ yang ramah, lalu ke $52$. Jawabannya adalah \omterm{def:g1:addition:def}{jumlah} seluruh
lompatan: $5 + 12 = 17$.}
\end{omfigure}

\section{Menaksir besarnya}

\begin{method}[Menaksir sebelum menghitung]\label{met:g3:addsub:estimate}
Sebelum menghitung secara bersusun, bulatkan bilangannya menjadi
bilangan yang ramah lalu hitunglah \emph{taksirannya}. Setelah
menghitung, bandingkanlah: hasil yang jauh dari taksiran menandakan
ada kekeliruan.
\end{method}

\begin{example}
Untuk $487 + 312$: taksirannya $500 + 300 = 800$. Hasil bersusunnya,
$799$, sudah dekat: masuk akal. Kalau kamu menemukan $1\,099$
(karena salah merapikan kolom), taksiran itu akan memperingatkanmu.
\end{example}

\section{Latihan}

\begin{exercise}[$\star$]\label{exo:g3:addsub:1}
Hitunglah secara bersusun: $346 + 253$; \quad $478 + 265$; \quad
$1\,536 + 2\,876$.
\end{exercise}

\begin{exercise}[$\star$]\label{exo:g3:addsub:2}
Hitunglah secara bersusun, lalu periksa dengan penjumlahan:
$587 - 234$; \quad $805 - 348$; \quad $1\,000 - 463$.
\end{exercise}

\begin{exercise}[$\star$]\label{exo:g3:addsub:3}
Hitunglah di kepala dengan membuat puluhan: $38 + 7$; \quad
$56 + 9$; \quad $45 + 28$ (tambahkan $30$, lalu ambil $2$).
\end{exercise}

\begin{exercise}[$\star$]\label{exo:g3:addsub:4}
Hitunglah di kepala dengan lompatan: $340 + 220$; \quad
$675 + 210$; \quad $512 - 300$; \quad $840 - 220$.
\end{exercise}

\begin{exercise}[$\star$]\label{exo:g3:addsub:5}
Hitunglah $63 - 47$ dengan membilang maju (seperti pada
\cref{ex:g3:addsub:missing}), lalu periksa secara bersusun.
\end{exercise}

\begin{exercise}[$\star$]\label{exo:g3:addsub:6}
Taksirlah dulu (bulatkan ke ratusan), lalu hitunglah dengan tepat:
$492 + 218$; \quad $913 - 388$.
\end{exercise}

\begin{exercise}[$\star$]\label{exo:g3:addsub:7}
Sebuah sekolah punya $348$ murid perempuan dan $296$ murid laki-laki.
Berapa murid seluruhnya? Tulislah operasinya, hitunglah, lalu jawablah dengan satu kalimat.
\end{exercise}

\begin{exercise}[$\star$]\label{exo:g3:addsub:8}
Sebuah buku setebal $224$ halaman; Nora sudah membaca $178$ halaman.
Berapa halaman yang tersisa untuk dibaca?
\end{exercise}

\begin{exercise}[$\star$]\label{exo:g3:addsub:9}
Lengkapilah bagian yang kosong (kerjakan kolom demi kolom):
\[
\begin{array}{r}
3\,?\,7 \\
+\ ?\,2\,? \\
\hline
8\,6\,2
\end{array}
\]
Apakah hanya ada satu cara mengisi ketiga lubang itu? Jelaskan.
\end{exercise}

\begin{exercise}[$\star\star$]\label{exo:g3:addsub:10}
Sebuah toko roti menjual $145$ croissant pada pagi hari dan $89$
pada sore hari; $17$ croissant tidak terjual sampai toko tutup.
Berapa croissant yang dipanggang hari itu? (Dua langkah.)
\end{exercise}

\begin{exercise}[$\star\star$]\label{exo:g3:addsub:11}
Teguh berkata: ``untuk menghitung $500 - 199$, aku menghitung
$500 - 200$ lalu menambahkan $1$.'' Benarkah dia? Hitunglah dengan kedua cara, lalu jelaskan triknya.

\end{exercise}
